\section{Contexto}\label{cap:contexto}

La Semana Santa es una de las fiestas culturales y religiosas más relevantes del calendario español. Detrás de esta tradición hay siglos de historia, ya que forma parte del patrimonio histórico y artístico del país. En ciudades como Sevilla, Málaga o Valladolid, entre otras, está declarada Fiesta de Interés Turístico Internacional. Esta celebración tiene un gran impacto tanto cultural como económico, debido a los miles de visitantes nacionales e internacionales que atrae cada año.

La estructura organizativa de la Semana Santa está formada por las cofradías o hermandades, las cuales son asociaciones religiosas encargadas de organizar procesiones y eventos durante estos días. Cada procesión sigue su itinerario por las calles de la ciudad, llevando consigo pasos de gran valor artístico que representan escenas de la Biblia. Esta organización de horarios, recorridos y actos requiere una gran coordinación por parte de las distintas cofradías y de la Junta de Cofradías de cada ciudad. Durante estos días, la ciudad experimenta imprevistos, como cortes de tráfico, cambios en la movilidad urbana y una gran cantidad de personas en puntos estratégicos del recorrido.

Desde el punto de vista social, la Semana Santa puede implicar dificultades de movilidad debido al desconocimiento de la ubicación de las distintas procesiones, además de los cambios de última hora en cuanto a los horarios de comienzo de cada procesión debido al mal tiempo o incluso a cancelaciones. Como esto se sale de la planificación inicial, los cofrades encuentran dificultades a la hora de comunicar estos cambios y decisiones de última hora a los espectadores y fieles. A día de hoy, son las distintas Juntas de Cofradías las que deben gestionar toda esta información de forma manual o con herramientas anticuadas.

En el contexto tecnológico actual, la mayoría de los turistas y ciudadanos hacen uso de sus teléfonos móviles como principal medio para consultar información. Aplicaciones como Google Maps \cite{googlemapsapi} permiten conocer en tiempo real la ubicación (geolocalización \cite{geolocalizacion}) de calles, monumentos, bares o incluso el estado del tráfico o del transporte público. Esta funcionalidad podría aplicarse a la Semana Santa, de forma que facilitaría a los visitantes y ciudadanos información constante y rápida sobre dónde se encuentra una procesión o si ha habido algún cambio o imprevisto en la misma.

Finalmente, debido a la cantidad de dinero que genera el turismo durante la Semana Santa, podemos concluir que es de gran interés para las ciudades que la experiencia del turista y del ciudadano sea positiva y que existan herramientas para que esto se haga realidad, de forma que una aplicación que mejore dicha experiencia repercutirá positivamente.
